\chapter{2016年上海卷（秋理）}

\section{填空题}

\begin{problem}
设 \(x\in\R\)，则不等式 \(|x-3|<1\) 的解集为\fillinblank{}.

\end{problem}

\begin{answer}
\((2,4)\)
\end{answer}

\begin{solution}
由 \(|x-3|<1\)，得 \(-1<x-3<1\)，所以 \(2<x<4\).
\end{solution}
\begin{problem}
设 \(z=\dfrac{3+2\i}{\i}\)，其中 \(\i\) 为虚数单位，则 \(\operatorname{Im}z=\fillinblank{}\).

\end{problem}

\begin{answer}
\(-3\)
\end{answer}

\begin{solution}
\(z=\frac{3+2\i}{\i}=2-3\i\)，所以 \(\operatorname{Im}z=-3\).
\end{solution}
\begin{problem}
已知平行直线 \(l_1:2x+y-1=0\)，\(l_2:2x+y+1=0\)，则 \(l_1,l_2\) 的距离是\fillinblank{}.

\end{problem}

\begin{answer}
\(\dfrac2{\sqrt5}\)
\end{answer}

\begin{solution}
两平行直线距离为 \(\frac{|(-1)-1|}{\sqrt{2^2+1^2}}=\frac2{\sqrt5}\).
\end{solution}
\begin{problem}
某次体检，\(6\) 位同学的身高（单位：米）分别为 \(1.72,1.78,1.75,1.80,1.69,1.77\)，则这组数据的中位数是\fillinblank{}（米）.

\end{problem}

\begin{answer}
\(1.76\)
\end{answer}

\begin{solution}
排序后为 \(1.69,1.72,1.75,1.77,1.78,1.80\)，中位数为 \(\frac{1.75+1.77}{2}=1.76\).
\end{solution}
\begin{problem}
已知点 \((3,9)\) 在函数 \(f(x)=1+a^x\) 的图像上，则 \(f(x)\) 的反函数 \(f^{-1}(x)=\fillinblank{}\).

\end{problem}

\begin{answer}
\(\log_2(x-1)\)
\end{answer}

\begin{solution}
由 \(9=1+a^3\) 得 \(a=2\). 故 \(f(x)=1+2^x\)，其反函数为 \(f^{-1}(x)=\log_2(x-1)\).
\end{solution}
\begin{problem}
如图，在正四棱柱 \(ABCD-A_1B_1C_1D_1\) 中，底面 \(ABCD\) 的边长为 \(3\)，\(BD_1\) 与底面所成角的大小为 \(\arctan\dfrac23\)，则该正四棱柱的高等于\fillinblank{}.

\bitmapfigure[max width=3.4cm,max totalheight=4.4cm]{img/2016/shanghai_science/q6_fig1.png}

\end{problem}

\begin{answer}
\(2\sqrt2\)
\end{answer}

\begin{solution}
设正四棱柱的高为 \(h\). 因为 \(BD=3\sqrt2\)，且 \(BD_1\) 与底面所成角满足 \(\tan\theta=\frac23\)，所以 \(\frac{h}{3\sqrt2}=\frac23\)，故 \(h=2\sqrt2\).
\end{solution}
\begin{problem}
方程 \(3\sin x=1+\cos2x\) 在区间 \([0,2\pi]\) 上的解为\fillinblank{}.

\end{problem}

\begin{answer}
\(\dfrac{\pi}{6},\dfrac{5\pi}{6}\)
\end{answer}

\begin{solution}
由 \(1+\cos2x=2-2\sin^2x\)，得 \(2\sin^2x+3\sin x-2=0\)，解得 \(\sin x=\frac12\) 或 \(\sin x=-2\). 舍去 \(\sin x=-2\)，故 \(x=\frac{\pi}{6}\) 或 \(x=\frac{5\pi}{6}\).
\end{solution}
\begin{problem}
在 \(\left(\sqrt[3]{x}-\dfrac2x\right)^n\) 的二项展开式中，所有项的二项式系数之和为 \(256\)，则常数项等于\fillinblank{}.

\end{problem}

\begin{answer}
\(112\)
\end{answer}

\begin{solution}
由二项式系数之和为 \(2^n=256\)，得 \(n=8\). 通项为 \(T_{k+1}=\mathrm C_8^k(\sqrt[3]{x})^{8-k}\left(-\frac2x\right)^k=\mathrm C_8^k(-2)^k x^{\frac{8-4k}{3}}\). 令 \(\frac{8-4k}{3}=0\)，得 \(k=2\). 常数项为 \(\mathrm C_8^2(-2)^2=112\).
\end{solution}
\begin{problem}
已知 \(\triangle ABC\) 的三边长分别为 \(3,5,7\)，则该三角形的外接圆半径等于\fillinblank{}.

\end{problem}

\begin{answer}
\(\dfrac{7\sqrt3}{3}\)
\end{answer}

\begin{solution}
由海伦公式，三角形面积为 \(\sqrt{\frac{15}{2}\cdot\frac92\cdot\frac52\cdot\frac12}=\frac{15\sqrt3}{4}\). 因此外接圆半径 \(R=\frac{abc}{4S}=\frac{3\cdot5\cdot7}{4\cdot\frac{15\sqrt3}{4}}=\frac{7\sqrt3}{3}\).
\end{solution}
\begin{problem}
设 \(a>0\)，\(b>0\)，若关于 \(x,y\) 的方程组
\[
\begin{cases}
ax+y=1,\\
x+by=1
\end{cases}
\]
无解，则 \(a+b\) 的取值范围是\fillinblank{}.

\end{problem}

\begin{answer}
\((2,+\infty)\)
\end{answer}

\begin{solution}
方程组无解需系数行列式为 \(ab-1=0\)，且两方程不重合. 因为 \(b=\frac1a\)，若两方程重合，则 \(a=1\)，\(b=1\). 因此 \(a>0\)，\(b=\frac1a\)，且 \(a\ne1\). 所以 \(a+b=a+\frac1a>2\)，取值范围为 \((2,+\infty)\).
\end{solution}
\begin{problem}
无穷数列 \(\{a_n\}\) 由 \(k\) 个不同的数组成，\(S_n\) 为 \(\{a_n\}\) 的前 \(n\) 项和. 若对任意的 \(n\in\N^*\)，\(S_n\in\{2,3\}\)，则 \(k\) 的最大值为\fillinblank{}.

\end{problem}

\begin{answer}
\(4\)
\end{answer}

\begin{solution}
因为 \(S_0=0\)，\(S_1\in\{2,3\}\)，故 \(a_1\) 只能为 \(2\) 或 \(3\). 当 \(n\ge2\) 时，\(a_n=S_n-S_{n-1}\)，其中 \(S_n,S_{n-1}\in\{2,3\}\)，所以 \(a_n\in\{-1,0,1\}\). 因此不同的项最多为 \(4\) 个. 例如 \(S_1,S_2,S_3,S_4,S_5,\ldots\) 取 \(3,2,3,3,2,\ldots\)，则 \(a_n\) 可出现 \(3,-1,1,0\) 四个不同的数，故最大值为 \(4\).
\end{solution}
\begin{problem}
在平面直角坐标系中，已知 \(A(1,0)\)，\(B(0,-1)\)，\(P\) 是曲线 \(y=\sqrt{1-x^2}\) 上一个动点，则 \(\overrightarrow{OP}\cdot\overrightarrow{BA}\) 的取值范围是\fillinblank{}.

\bitmapfigure[max width=3.3cm,max totalheight=3.3cm]{img/2016/shanghai_science/q12_fig1.png}

\end{problem}

\begin{answer}
\([-1,\sqrt2]\)
\end{answer}

\begin{solution}
设 \(P(x,y)\)，则 \(x^2+y^2=1\)，且 \(y\ge0\). 因为 \(\overrightarrow{OP}=(x,y)\)，\(\overrightarrow{BA}=(1,1)\)，所以 \(\overrightarrow{OP}\cdot\overrightarrow{BA}=x+y\). 在上半圆上，\(x+y\) 的最大值为 \(\sqrt2\)，最小值在点 \((-1,0)\) 处取得，为 \(-1\).
\end{solution}
\begin{problem}
设 \(a,b\in\R\)，\(c\in[0,2\pi)\)，若对任意实数 \(x\) 都有 \(2\sin\!\left(3x-\dfrac{\pi}{3}\right)=a\sin(bx+c)\)，则满足条件的有序实数组 \((a,b,c)\) 的组数为\fillinblank{}.

\end{problem}

\begin{answer}
\(4\)
\end{answer}

\begin{solution}
两边为同一正弦函数，需 \(|a|=2\)，且 \(|b|=3\). 当 \(b=3\) 时，可取 \((a,c)=(2,\frac{5\pi}{3})\) 或 \((-2,\frac{2\pi}{3})\). 当 \(b=-3\) 时，可取 \((a,c)=(-2,\frac{\pi}{3})\) 或 \((2,\frac{4\pi}{3})\). 因此共有 \(4\) 组.
\end{solution}
\begin{problem}
如图，在平面直角坐标系 \(xOy\) 中，\(O\) 为正八边形 \(A_1A_2\cdots A_8\) 的中心，\(A_1(1,0)\)，任取不同的两点 \(A_i,A_j\)，点 \(P\) 满足 \(\overrightarrow{OP}+\overrightarrow{OA_i}+\overrightarrow{OA_j}=\bs0\)，则点 \(P\) 落在第一象限的概率是\fillinblank{}.

\bitmapfigure[max width=4cm,max totalheight=4cm]{img/2016/shanghai_science/q14_fig1.png}

\end{problem}

\begin{answer}
\(\dfrac5{28}\)
\end{answer}

\begin{solution}
共有 \(\mathrm C_8^2=28\) 种不同取法. 因为 \(\overrightarrow{OP}=-(\overrightarrow{OA_i}+\overrightarrow{OA_j})\)，所以 \(P\) 在第一象限等价于 \(\overrightarrow{OA_i}+\overrightarrow{OA_j}\) 在第三象限. 逐对计数可得符合条件的点对为 \((A_4,A_7)\)，\((A_5,A_6)\)，\((A_5,A_7)\)，\((A_5,A_8)\)，\((A_6,A_7)\)，共 \(5\) 对. 故概率为 \(\frac5{28}\).
\end{solution}
\section{单选题}

\begin{problem}
设 \(a\in\R\)，则“\(a>1\)”是“\(a^2>1\)”的（\quad）

\choices
    {充分非必要条件}
    {必要非充分条件}
    {充要条件}
    {既非充分也非必要条件}

\end{problem}

\begin{answer}
A
\end{answer}

\begin{solution}
由 \(a>1\) 可推出 \(a^2>1\). 但 \(a^2>1\) 可得 \(a>1\) 或 \(a<-1\)，不能推出 \(a>1\). 故选 \(A\).
\end{solution}
\begin{problem}
下列极坐标方程中，对应的曲线为下图的是（\quad）

\bitmapfigure[max width=4.2cm,max totalheight=3.4cm]{img/2016/shanghai_science/q16_fig1.png}

\choices
    {\(\rho=6+5\cos\theta\)}
    {\(\rho=6+5\sin\theta\)}
    {\(\rho=6-5\cos\theta\)}
    {\(\rho=6-5\sin\theta\)}

\end{problem}

\begin{answer}
C
\end{answer}

\begin{solution}
图形关于 \(x\) 轴对称，且在正 \(x\) 轴方向半径最小，在负 \(x\) 轴方向半径最大. 因此应满足 \(\theta=0\) 时 \(\rho\) 最小，\(\theta=\pi\) 时 \(\rho\) 最大，对应 \(\rho=6-5\cos\theta\). 故选 \(C\).
\end{solution}
\begin{problem}
已知无穷等比数列 \(\{a_n\}\) 的公比为 \(q\)，前 \(n\) 项和为 \(S_n\)，且 \(\lim_{n\to\infty}S_n=S\)，下列条件中，使得 \(2S_n<S\)（\(n\in\N^*\)）恒成立的是（\quad）

\choices
    {\(a_1>0,\ 0.6<q<0.7\)}
    {\(a_1<0,\ -0.7<q<-0.6\)}
    {\(a_1>0,\ 0.7<q<0.8\)}
    {\(a_1<0,\ -0.8<q<-0.7\)}

\end{problem}

\begin{answer}
B
\end{answer}

\begin{solution}
由 \(|q|<1\)，\(S=\frac{a_1}{1-q}\)，且 \(S_n=S(1-q^n)\). 条件 \(2S_n<S\) 等价于 \(S(1-2q^n)<0\). 若 \(a_1<0\)，则 \(S<0\)，需 \(1-2q^n>0\) 对任意 \(n\) 成立. 当 \(-0.7<q<-0.6\) 时，奇数次 \(q^n<0\)，偶数次 \(q^n\le q^2<0.49\)，故恒有 \(1-2q^n>0\). 故选 \(B\).
\end{solution}
\begin{problem}
设 \(f(x),g(x),h(x)\) 是定义域为 \(\R\) 的三个函数. 对于命题：\(\circled{1}\) 若 \(f(x)+g(x)\)，\(f(x)+h(x)\)，\(g(x)+h(x)\) 均为增函数，则 \(f(x),g(x),h(x)\) 中至少有一个为增函数；\(\circled{2}\) 若 \(f(x)+g(x)\)，\(f(x)+h(x)\)，\(g(x)+h(x)\) 均是以 \(T\) 为周期的函数，则 \(f(x),g(x),h(x)\) 均是以 \(T\) 为周期的函数. 下列判断正确的是（\quad）

\choices
    {\(\circled{1}\) 和 \(\circled{2}\) 均为真命题}
    {\(\circled{1}\) 和 \(\circled{2}\) 均为假命题}
    {\(\circled{1}\) 为真命题，\(\circled{2}\) 为假命题}
    {\(\circled{1}\) 为假命题，\(\circled{2}\) 为真命题}

\end{problem}

\begin{answer}
D
\end{answer}

\begin{solution}
\(\circled{1}\) 为假命题. 对任意 \(x_1<x_2\)，设三个函数增量分别为 \(\Delta f,\Delta g,\Delta h\). 即使 \(\Delta f+\Delta g>0\)，\(\Delta f+\Delta h>0\)，\(\Delta g+\Delta h>0\)，也可能有 \(\Delta f<0\). 例如取 \(\Delta f=-1\)，\(\Delta g=2\)，\(\Delta h=2\). \(\circled{2}\) 为真命题，因为三组和函数均以 \(T\) 为周期，由 \(f+g\) 与 \(f+h\) 的周期性相减可得 \(g(x+T)-h(x+T)=g(x)-h(x)\)，再结合 \(g+h\) 的周期性可得 \(g(x+T)=g(x)\). 继续结合 \(f+g\)、\(g+h\) 的周期性，可得 \(f,h\) 也均以 \(T\) 为周期. 故选 \(D\).
\end{solution}
\section{解答题}

\begin{problem}

将边长为 \(1\) 的正方形 \(AA_1O_1O\)（及其内部）绕 \(OO_1\) 旋转一周形成圆柱，如图，\(\widehat{AC}\) 长为 \(\dfrac{2\pi}{3}\)，\(\widehat{A_1B_1}\) 长为 \(\dfrac{\pi}{3}\)，其中 \(B_1\) 与 \(C\) 在平面 \(AA_1O_1O\) 的同侧.

\begin{enumerate}
\item 求三棱锥 \(C-O_1A_1B_1\) 的体积；
\item 求异面直线 \(B_1C\) 与 \(AA_1\) 所成角的大小.
\end{enumerate}

\bitmapfigure[max width=4.5cm,max totalheight=4cm]{img/2016/shanghai_science/q19_fig1.png}

\end{problem}

\begin{solution}
建立空间直角坐标系，使 \(O(0,0,0)\)，\(O_1(0,0,1)\)，\(A(1,0,0)\)，\(A_1(1,0,1)\). 因为底面圆半径为 \(1\)，\(\widehat{AC}=\frac{2\pi}{3}\)，\(\widehat{A_1B_1}=\frac{\pi}{3}\)，且 \(B_1\) 与 \(C\) 在同侧，所以可取 \(C\left(-\frac12,\frac{\sqrt3}{2},0\right)\)，\(B_1\left(\frac12,\frac{\sqrt3}{2},1\right)\).

（1）在平面 \(z=1\) 内，\(\triangle O_1A_1B_1\) 的面积为 \(\frac12\cdot1\cdot1\cdot\sin\frac{\pi}{3}=\frac{\sqrt3}{4}\). 点 \(C\) 到平面 \(O_1A_1B_1\) 的距离为 \(1\)，所以三棱锥体积为 \(V=\frac13\cdot\frac{\sqrt3}{4}\cdot1=\frac{\sqrt3}{12}\).

（2）\(\overrightarrow{B_1C}=(-1,0,-1)\)，\(\overrightarrow{AA_1}=(0,0,1)\). 设异面直线所成角为 \(\theta\)，则 \(\cos\theta=\frac{|\overrightarrow{B_1C}\cdot\overrightarrow{AA_1}|}{|\overrightarrow{B_1C}||\overrightarrow{AA_1}|}=\frac1{\sqrt2}\)，故 \(\theta=\frac{\pi}{4}\).
\end{solution}
\begin{problem}

有一块正方形菜地 \(EFGH\)，\(EH\) 所在直线是一条小河，收获的蔬菜可送到 \(F\) 点或河边运走. 于是，菜地分为两个区域 \(S_1\) 和 \(S_2\)，其中 \(S_1\) 中的蔬菜运到河边较近，\(S_2\) 中的蔬菜运到 \(F\) 点较近，而菜地内 \(S_1\) 和 \(S_2\) 的分界线 \(C\) 上的点到河边与到 \(F\) 点的距离相等. 现建立平面直角坐标系，其中原点 \(O\) 为 \(EF\) 的中点，点 \(F\) 的坐标为 \((1,0)\)，如图.

\begin{enumerate}
\item 求菜地内的分界线 \(C\) 的方程；
\item 菜农从蔬菜运量估计出 \(S_1\) 面积是 \(S_2\) 面积的两倍，由此得到 \(S_1\) 面积的“经验值”为 \(\dfrac83\). 设 \(M\) 是 \(C\) 上纵坐标为 \(1\) 的点，请计算以 \(EH\) 为一边、另有一边过点 \(M\) 的矩形的面积，及五边形 \(EOMGH\) 的面积，并判别哪一个更接近于 \(S_1\) 面积的“经验值”.
\end{enumerate}

\bitmapfigure[max width=4.4cm,max totalheight=4.2cm]{img/2016/shanghai_science/q20_fig1.png}

\end{problem}

\begin{solution}
（1）设分界线 \(C\) 上一点为 \((x,y)\). 因为点到河边 \(EH:x=-1\) 的距离为 \(x+1\)，点到 \(F(1,0)\) 的距离为 \(\sqrt{(x-1)^2+y^2}\)，所以 \(x+1=\sqrt{(x-1)^2+y^2}\). 化简得 \(y^2=4x\). 因此菜地内分界线 \(C\) 的方程为 \(y^2=4x\)（\(0\le y\le2\)）.

（2）由 \(M\) 的纵坐标为 \(1\)，得 \(1=4x_M\)，所以 \(M\left(\frac14,1\right)\). 以 \(EH\) 为一边、另一边过点 \(M\) 的矩形高为 \(2\)，宽为 \(\frac14-(-1)=\frac54\)，面积为 \(\frac52\). 五边形 \(EOMGH\) 的顶点可取 \(E(-1,0)\)，\(O(0,0)\)，\(M(\frac14,1)\)，\(G(1,2)\)，\(H(-1,2)\)，用割补或鞋带公式得面积为 \(\frac{11}{4}\). 因为 \(\left|\frac52-\frac83\right|=\frac16\)，\(\left|\frac{11}{4}-\frac83\right|=\frac1{12}\)，所以五边形 \(EOMGH\) 的面积更接近 \(S_1\) 面积的“经验值”.
\end{solution}
\begin{problem}

双曲线 \(x^2-\dfrac{y^2}{b^2}=1\)（\(b>0\)）的左、右焦点分别为 \(F_1,F_2\)，直线 \(l\) 过 \(F_2\) 且与双曲线交于 \(A,B\) 两点.

\begin{enumerate}
\item 若 \(l\) 的倾斜角为 \(\dfrac{\pi}{2}\)，\(\triangle F_1AB\) 是等边三角形，求双曲线的渐近线方程；
\item 设 \(b=\sqrt3\)，若 \(l\) 的斜率存在，且 \((\overrightarrow{F_1A}+\overrightarrow{F_1B})\cdot\overrightarrow{AB}=0\)，求 \(l\) 的斜率.
\end{enumerate}

\end{problem}

\begin{solution}
（1）设 \(c=\sqrt{1+b^2}\)，则 \(F_1(-c,0)\)，\(F_2(c,0)\). 因为 \(l\) 的倾斜角为 \(\frac{\pi}{2}\)，所以 \(l:x=c\). 代入双曲线得 \(y=\pm b^2\)，于是 \(A(c,b^2)\)，\(B(c,-b^2)\). 因为 \(\triangle F_1AB\) 是等边三角形，所以 \(\sqrt{(2c)^2+b^4}=2b^2\). 即 \(4(1+b^2)+b^4=4b^4\)，解得 \(b^2=2\). 故渐近线方程为 \(y=\pm\sqrt2x\).

（2）当 \(b=\sqrt3\) 时，\(F_1(-2,0)\)，\(F_2(2,0)\). 设 \(l:y=k(x-2)\)，与双曲线交点横坐标为 \(x_1,x_2\). 联立得 \(\left(1-\frac{k^2}{3}\right)x^2+\frac{4k^2}{3}x-\frac{4k^2}{3}-1=0\)，故 \(x_1+x_2=-\frac{4k^2}{3-k^2}\). 由 \((\overrightarrow{F_1A}+\overrightarrow{F_1B})\cdot\overrightarrow{AB}=0\)，可得 \((1+k^2)(x_1+x_2)+4(1-k^2)=0\). 代入并整理得 \(k^2=\frac35\)，所以 \(l\) 的斜率为 \(\pm\sqrt{\frac35}\).
\end{solution}
\begin{problem}

已知 \(a\in\R\)，函数 \(f(x)=\log_2\!\left(\dfrac1x+a\right)\).

\begin{enumerate}
\item 当 \(a=5\) 时，解不等式 \(f(x)>0\)；
\item 若关于 \(x\) 的方程 \(f(x)-\log_2[(a-4)x+2a-5]=0\) 的解集中恰有一个元素，求 \(a\) 的取值范围；
\item 设 \(a>0\)，若对任意 \(t\in\left[\dfrac12,1\right]\)，函数 \(f(x)\) 在区间 \([t,t+1]\) 上的最大值和最小值的差不超过 \(1\)，求 \(a\) 的取值范围.
\end{enumerate}

\end{problem}

\begin{solution}
（1）当 \(a=5\) 时，定义域满足 \(\frac1x+5>0\)，即 \(x\in(-\infty,-\frac15)\cup(0,+\infty)\). 不等式 \(f(x)>0\) 等价于 \(\frac1x+5>1\)，即 \(\frac1x>-4\). 结合定义域，得解集为 \((-\infty,-\frac14)\cup(0,+\infty)\).

（2）由对数相等，得 \(\frac1x+a=(a-4)x+2a-5\). 因为 \(x\ne0\)，两边乘以 \(x\)，得 \((a-4)x^2+(a-5)x-1=0\). 当 \(a\ne4\) 时，方程可化为 \((x+1)((a-4)x-1)=0\)，即 \(x=-1\) 或 \(x=\frac1{a-4}\). 其中 \(x=-1\) 是原方程解当且仅当 \(a>1\)，\(x=\frac1{a-4}\) 是原方程解当且仅当 \(a>2\). 因此当 \(1<a\le2\) 时，解集中恰有一个元素；当 \(a=3\) 时两根重合为 \(-1\)，解集中也恰有一个元素. 当 \(a=4\) 时，原代数方程化为 \(-x-1=0\)，得唯一解 \(x=-1\)，且满足定义域. 故 \(a\) 的取值范围为 \((1,2]\cup\{3,4\}\).

（3）当 \(a>0\) 时，\(f(x)\) 在 \((0,+\infty)\) 上单调递减. 因此 \(f(x)\) 在 \([t,t+1]\) 上的最大值和最小值之差为 \(f(t)-f(t+1)\). 由 \(f(t)-f(t+1)\le1\)，得 \(\log_2\frac{\frac1t+a}{\frac1{t+1}+a}\le1\)，即 \(\frac{1-t}{t(t+1)}\le a\). 函数 \(h(t)=\frac{1-t}{t(t+1)}\) 在 \([\frac12,1]\) 上单调递减，最大值为 \(h(\frac12)=\frac23\)，故 \(a\ge\frac23\).
\end{solution}
\begin{problem}

若无穷数列 \(\{a_n\}\) 满足：只要 \(a_p=a_q\)（\(p,q\in\N^*\)），必有 \(a_{p+1}=a_{q+1}\)，则称 \(\{a_n\}\) 具有性质 \(P\).

\begin{enumerate}
\item 若 \(\{a_n\}\) 具有性质 \(P\)，且 \(a_1=1\)，\(a_2=2\)，\(a_4=3\)，\(a_5=2\)，\(a_6+a_7+a_8=21\)，求 \(a_3\)；
\item 若无穷数列 \(\{b_n\}\) 是等差数列，无穷数列 \(\{c_n\}\) 是公比为正数的等比数列，\(b_1=c_5=1\)，\(b_5=c_1=81\)，\(a_n=b_n+c_n\)，判断 \(\{a_n\}\) 是否具有性质 \(P\)，并说明理由；
\item 设 \(\{b_n\}\) 是无穷数列，已知 \(a_{n+1}=b_n+\sin a_n\)（\(n\in\N^*\)），求证：“对任意 \(a_1\)，\(\{a_n\}\) 都具有性质 \(P\)”的充要条件为“\(\{b_n\}\) 是常数列”.
\end{enumerate}

\end{problem}

\begin{solution}
（1）因为 \(a_2=a_5=2\)，由性质 \(P\) 得 \(a_3=a_6\). 又由 \(a_3=a_6\)，得 \(a_4=a_7\)，所以 \(a_7=3\). 再由 \(a_4=a_7=3\)，得 \(a_5=a_8\)，所以 \(a_8=2\). 因此 \(a_6+a_7+a_8=a_3+3+2=21\)，故 \(a_3=16\).

（2）由 \(b_1=1\)，\(b_5=81\)，得等差数列 \(\{b_n\}\) 的公差为 \(20\)，所以 \(b_n=20n-19\). 由 \(c_1=81\)，\(c_5=1\)，且公比为正，得公比为 \(\frac13\)，所以 \(c_n=81\left(\frac13\right)^{n-1}\). 因此 \(a_1=82\)，\(a_5=82\)，但 \(a_2=48\)，\(a_6=\frac{304}{3}\)，所以由 \(a_1=a_5\) 不能推出 \(a_2=a_6\). 故 \(\{a_n\}\) 不具有性质 \(P\).

（3）充分性：若 \(\{b_n\}\) 是常数列，设 \(b_n=b\). 若 \(a_p=a_q\)，则 \(a_{p+1}=b+\sin a_p=b+\sin a_q=a_{q+1}\)，所以 \(\{a_n\}\) 具有性质 \(P\).

必要性：若对任意 \(a_1\)，\(\{a_n\}\) 都具有性质 \(P\). 对任意 \(m\ge2\)，由于 \(a_m\) 是关于 \(a_1\) 的连续函数，且当 \(a_1\to+\infty\) 时，\(a_m-a_1\to-\infty\)，当 \(a_1\to-\infty\) 时，\(a_m-a_1\to+\infty\)，故存在 \(a_1\) 使 \(a_m=a_1\). 由性质 \(P\) 得 \(a_{m+1}=a_2\). 又 \(a_{m+1}=b_m+\sin a_m\)，\(a_2=b_1+\sin a_1\)，且 \(a_m=a_1\)，所以 \(b_m=b_1\). 因为 \(m\ge2\) 任意，所以 \(\{b_n\}\) 是常数列. 综上，命题得证.
\end{solution}
